\begin{abstractcn}
本报告围绕虚拟现实课程的作业展开。首先，针对我小组课堂汇报中的人脸分析与识别主题，完成了浏览器端人脸分析项目的详细讲解和 dlib 人脸识别示例的解析。基于 face-api.js，分析了静态网页中模型加载、图片输入、摄像头视频流处理与 Canvas 可视化流程，实现了人脸检测、关键点定位、表情识别以及年龄性别估计；结合 dlib 官方度量学习示例，进一步剖析了训练数据组织、mini-batch 采样、ResNet 描述子网络、loss\_metric 损失函数与距离阈值验证机制，阐明人脸描述子模型如何通过特征距离完成身份比较。此外，在完成课程要求的基础上，本文还开展了面向三维虚拟场景搭建的附加实践：使用 Unity VR 模板构建森林场景，导入森林、花卉与动物资源，修复旧资源在 Unity 6 + URP 环境下的粉色材质问题，接入 XR Origin 与移动系统，并完成 Windows 平台的打包运行。报告涵盖了模型推理、工程实现、资源兼容与场景发布等课程实践内容。

\keywords{人脸分析；人脸识别；face-api.js；dlib；度量学习；Unity；虚拟现实；URP}
\end{abstractcn}

\begin{abstracten}
This report summarizes the coursework completed for the Virtual Reality course. First, based on the topic of face analysis and recognition presented by my group in class, it provides a detailed explanation of a browser-based face analysis project and analyzes a dlib face recognition example. Using face-api.js, the report examines the workflow of model loading, image input, camera video stream processing, and Canvas-based visualization in a static web page, implementing face detection, facial landmark localization, expression recognition, and age and gender estimation. It also discusses dlib's official metric learning example, including training data organization, mini-batch sampling, the ResNet descriptor network, the loss\_metric loss function, and distance-threshold verification, explaining how a face descriptor model compares identities through feature distances. In addition to the required coursework, this report includes an extended practice project on constructing a three-dimensional virtual scene. A forest scene was built with the Unity VR template, forest, flower, and animal assets were imported, pink material issues caused by legacy assets in Unity 6 with URP were resolved, XR Origin and locomotion systems were integrated, and the project was packaged for Windows. Overall, the report covers practical work in model inference, engineering implementation, asset compatibility, and scene deployment.

\enkeywords{face analysis; face recognition; face-api.js; dlib; metric learning; Unity; virtual reality; URP}
\end{abstracten}
